
\subsection{Variational Autoencoders (VAE)}
\paragraph{Latent-variable model and variational inference.}
We consider a latent-variable generative model
\[
p_\theta(x,z)=p_\theta(x\mid z)\,p(z),
\]
where $x\in\mathbb{R}^{D}$ denotes an observed image and $z\in\mathbb{R}^{d}$ is a continuous latent variable. Since the exact posterior $p_\theta(z\mid x)$ is typically intractable, we introduce an amortized variational approximation (encoder)
\[
q_\phi(z\mid x).
\]
For Dynamic MNIST, we use a Bernoulli observation model
\[
p_\theta(x\mid z)=\prod_{i=1}^D \mathrm{Bernoulli}\bigl(x_i; \pi_\theta(z)_i\bigr),
\]
where $\pi_\theta(z)\in(0,1)^D$ is the decoder output after a sigmoid nonlinearity.

\paragraph{ELBO objective and KL regularization.}
Training proceeds by maximizing the evidence lower bound (ELBO) on the log-likelihood $\log p_\theta(x)$ \citep{kingma2014vae,rezende2014stochastic}:
\[
\log p_\theta(x)\;\ge\;\mathcal{L}_{\mathrm{ELBO}}(x)
=
\mathbb{E}_{q_\phi(z\mid x)}\left[\log p_\theta(x\mid z)\right]
-
\mathrm{KL}\!\left(q_\phi(z\mid x)\,\|\,p(z)\right).
\]
The first term is a reconstruction (negative cross-entropy for Bernoulli pixels), while the KL divergence regularizes the approximate posterior toward the prior, controlling latent capacity and enabling sampling from $p(z)$ at test time.

\paragraph{Reparameterization trick.}
We parameterize $q_\phi(z\mid x)$ as a diagonal Gaussian,
\[
q_\phi(z\mid x)=\mathcal{N}\!\left(z;\mu_\phi(x),\mathrm{diag}(\sigma_\phi^2(x))\right),
\]
and apply the reparameterization trick to enable low-variance gradient estimation:
\[
z=\mu_\phi(x)+\sigma_\phi(x)\odot\epsilon,\qquad \epsilon\sim\mathcal{N}(0,I).
\]

\paragraph{Enhanced-capacity objective: flexible prior via pseudo-inputs.}
A core limitation of standard VAEs is that a simple fixed prior (e.g., $\mathcal{N}(0,I)$) can over-regularize $q_\phi(z\mid x)$, encouraging inactive latent dimensions and contributing to posterior collapse, particularly when the decoder is expressive. Following the referenced method, we replace the fixed prior with a learnable, data-adaptive prior defined through \emph{pseudo-inputs} $\{u_k\}_{k=1}^K$:
\[
p_{\lambda}(z)
=
\frac{1}{K}\sum_{k=1}^{K} q_\phi(z\mid u_k),
\]
where $u_k$ are optimized parameters in input space (or are generated by a small network) and $q_\phi(z\mid u_k)$ reuses the encoder. This prior increases representational capacity by forming a multimodal latent distribution aligned with the aggregate posterior, thereby relaxing an overly restrictive KL constraint. The training objective becomes:
\[
\mathcal{L}(x)=
\mathbb{E}_{q_\phi(z\mid x)}\left[\log p_\theta(x\mid z)\right]
-
\mathrm{KL}\!\left(q_\phi(z\mid x)\,\|\,p_{\lambda}(z)\right),
\]
with the KL now computed against the learned mixture prior. In practice, this modification can improve latent utilization and sample quality by avoiding a mismatch between $q_\phi(z\mid x)$ and a simplistic prior.

\paragraph{Hierarchical latent architecture.}
To further increase capacity, we implement a hierarchical latent model with two stochastic layers (e.g., $z_2\rightarrow z_1\rightarrow x$), where the generative model factorizes as
\[
p_\theta(x,z_1,z_2)=p_\theta(x\mid z_1)\,p_\theta(z_1\mid z_2)\,p(z_2),
\]
and the inference model factorizes as
\[
q_\phi(z_1,z_2\mid x)=q_\phi(z_2\mid x)\,q_\phi(z_1\mid x,z_2).
\]
The corresponding ELBO is
\[
\mathcal{L}_{\mathrm{hier}}(x)=
\mathbb{E}_{q_\phi(z_1,z_2\mid x)}\left[\log p_\theta(x\mid z_1)\right]
-
\mathbb{E}_{q_\phi(z_2\mid x)}\left[
\mathrm{KL}\!\left(q_\phi(z_1\mid x,z_2)\,\|\,p_\theta(z_1\mid z_2)\right)\right]
-
\mathrm{KL}\!\left(q_\phi(z_2\mid x)\,\|\,p(z_2)\right),
\]
where the top-level prior $p(z_2)$ may be chosen as a standard normal or replaced by the pseudo-input prior $p_\lambda(z_2)$, depending on the assignment’s specification.

\paragraph{Likelihood estimation.}
Since $\log p_\theta(x)=\log\int p_\theta(x\mid z)p(z)\,dz$ is not directly tractable, we estimate log-likelihood using an importance-weighted (IWAE-style) estimator \citep{burda2016iwae}. For $K$ importance samples $z^{(k)}\sim q_\phi(z\mid x)$,
\[
\log p_\theta(x)\approx
\log\left(\frac{1}{K}\sum_{k=1}^{K}
\frac{p_\theta(x\mid z^{(k)})\,p(z^{(k)})}{q_\phi(z^{(k)}\mid x)}\right).
\]

\paragraph{Diagnostics for posterior behavior and collapse.}
To analyze latent utilization, we report (i) KL contributions per dimension (or per group in hierarchical models), (ii) active units (e.g., dimensions with $\mathrm{Var}_x[\mathbb{E}(z\mid x)]$ above a threshold), and (iii) aggregate posterior statistics compared with the prior. We complement these with latent traversals and PCA of inferred codes to test whether principal directions correspond to meaningful factors of variation.


\subsection{Conditional GAN for Data Augmentation (WGAN-GP)}
\paragraph{Problem setting and data splits.}
We consider FashionMNIST under a limited-data regime. The training set is partitioned into two \emph{non-overlapping} subsets: one for GAN training and one for classifier training. The GAN is used to generate additional labeled samples for augmentation; the downstream classifier is trained on either real-only or real+synthetic data and evaluated on a held-out test set.

\paragraph{Conditional generator.}
The generator is conditional on class label $y\in\{1,\dots,C\}$ and noise $z\sim\mathcal{N}(0,I)$. A learnable embedding $e_G(y)\in\mathbb{R}^{d_y}$ is concatenated with $z$ to form the generator input:
\[
\tilde{z} = [z;\,e_G(y)],\qquad x_{\mathrm{fake}} = G_\psi(\tilde{z},y).
\]
The architecture follows the assignment’s specified ConvTranspose/upsampling blocks, including attention modules at designated resolutions.

\paragraph{Projection discriminator (conditional critic).}
The discriminator (critic) follows the projection discriminator formulation \citep{miyato2018projection}. Let $\phi_\omega(x)\in\mathbb{R}^{d_f}$ denote the discriminator feature vector before the final output layer, and let $e_D(y)\in\mathbb{R}^{d_f}$ be the label embedding in the discriminator. The conditional score is computed as
\[
D_\omega(x,y)=
s_\omega(\phi_\omega(x)) + \langle \phi_\omega(x),\, e_D(y)\rangle,
\]
where $s_\omega(\cdot)$ is a learned scalar head and $\langle\cdot,\cdot\rangle$ denotes an inner product. This construction enforces class-conditional discrimination without explicitly concatenating labels to image channels.

\paragraph{Spectral normalization.}
To stabilize adversarial training, we apply spectral normalization (SN) to convolutional (and transposed convolutional) layers \citep{miyato2018sn}. SN constrains each layer’s operator norm by normalizing its weight matrix $W$ via its largest singular value $\sigma_{\max}(W)$:
\[
\bar{W} = \frac{W}{\sigma_{\max}(W)},
\]
which controls the Lipschitz constant of the discriminator and improves training stability.

\paragraph{Attention mechanisms.}
We incorporate two attention modules, following the assignment definitions.

\emph{Self-attention} (as in SAGAN \citep{zhang2019sagan}) captures long-range spatial dependencies. Given feature maps $x\in\mathbb{R}^{B\times C\times H\times W}$ with $N=HW$, we compute projections $f,g,h$ and attention weights $\beta\in\mathbb{R}^{B\times N\times N}$:
\[
\beta = \mathrm{softmax}(f^\top g),
\]
and produce an output $o$ aggregated from $h$ (and then mapped through $v$), followed by a residual connection with a learnable scalar $\gamma$:
\[
y = \gamma o + x.
\]

\emph{Channel attention} (SE-style \citep{hu2018senet}) recalibrates channel-wise responses. Global average pooling produces a channel descriptor $c\in\mathbb{R}^{B\times C}$, which is passed through two fully-connected layers with ReLU and a sigmoid gate to generate channel weights that modulate the original feature map.

\paragraph{WGAN-GP objective with regularization.}
We use the Wasserstein GAN objective with gradient penalty (WGAN-GP) \citep{gulrajani2017improved}. The critic aims to maximize the Wasserstein estimate:
\[
\mathcal{W} \approx \mathbb{E}[D(x_{\mathrm{real}},y)] - \mathbb{E}[D(x_{\mathrm{fake}},y)],
\]
subject to a soft 1-Lipschitz constraint enforced by the gradient penalty:
\[
\mathcal{L}_{\mathrm{GP}}=
\lambda_{\mathrm{GP}}
\mathbb{E}_{\hat{x}}\left(\left\|\nabla_{\hat{x}} D(\hat{x},y)\right\|_2 - 1\right)^2,
\qquad \lambda_{\mathrm{GP}}=10,
\]
where $\hat{x}=\epsilon x_{\mathrm{real}}+(1-\epsilon)x_{\mathrm{fake}}$ and $\epsilon\sim\mathcal{U}(0,1)$.

In addition, we incorporate:
(i) a drift penalty to prevent critic outputs from drifting:
\[
\mathcal{L}_{\mathrm{drift}}=\epsilon_{\mathrm{drift}}\;\mathbb{E}\left[D(x_{\mathrm{real}},y)^2\right],\qquad \epsilon_{\mathrm{drift}}=0.001,
\]
and (ii) discriminator embedding regularization:
\[
\mathcal{L}_{\mathrm{emb}}=\lambda_{\mathrm{emb}}\;\|e_D\|_2^2,\qquad \lambda_{\mathrm{emb}}=0.001.
\]
The discriminator loss (to minimize) is therefore:
\[
\mathcal{L}_D=
-\Bigl(\mathbb{E}[D(x_{\mathrm{real}},y)]-\mathbb{E}[D(x_{\mathrm{fake}},y)]\Bigr)
+\mathcal{L}_{\mathrm{GP}}+\mathcal{L}_{\mathrm{drift}}+\mathcal{L}_{\mathrm{emb}},
\]
and the generator loss is:
\[
\mathcal{L}_G=-\mathbb{E}[D(x_{\mathrm{fake}},y)].
\]

\paragraph{Optimization and training protocol.}
We use Adam optimizers with the assignment’s hyperparameters: learning rates $\mathrm{LR}(G)=2\times10^{-4}$ and $\mathrm{LR}(D)=4\times10^{-4}$, $\beta_1=0.0$, $\beta_2=0.9$, and an exponential decay scheduler with $\gamma=0.95$ applied every 100 steps. Weight decay is grouped such that embedding parameters receive decay $10^{-3}$ and other parameters receive zero decay, matching the specification. We update the discriminator $N_{\mathrm{critic}}=5$ times per generator update, detach fake samples during critic updates, and log the Wasserstein estimate, gradient penalty, drift, embedding regularization, and total losses at each step.

\paragraph{Sampling protocol and augmentation evaluation.}
To monitor training progression, we generate class-conditional grids at fixed steps using fixed noise and fixed labels (covering all classes), enabling controlled qualitative comparisons across time. For augmentation, we generate a specified number of synthetic samples per class and train a classifier under two regimes: real-only and real+synthetic. We then compare overall accuracy, class-wise performance, and error patterns to quantify whether generative augmentation improves generalization in the limited data setting.
